% Copyright 2026 The terCAD team. All rights reserved.
% Use of this content is governed by a CC BY-NC-ND 4.0 license that can be found in the LICENSE file.

\subsection{Architecting Project} \label{architecture}
\markboth{Prototyping}{Architecting Project}

When a mobile project begins under uncertainty, it is tempting to stay in an "architecture-free"
mode for as long as possible: wire up a first route, get a screen rendering inside the emulator,
and iterate from there. That works for a single proof of concept. Once the native shell is
running (\cref{native-components}), however, a NativePHP codebase quickly reveals a hidden
tension: two execution worlds---a \textbf{native mobile shell} and an \textbf{embedded Laravel
web runtime}---share the same repository and the same PHP process. Without clear boundaries,
platform workarounds, storage details, and UI logic rapidly tangle into code that is painful to
debug and impossible to test without a phone.


\subsubsection{Separating Layers and Responsibilities}

Just as in any cross-platform application, we want a clean boundary between what renders the
interface, what handles requests, and what manages state. Rather than inventing a rigid folder
hierarchy that fights Laravel's conventions, we separate \emph{logical layers} conceptually and
keep files organized around their framework roles and feature targets:

\hspace{0.5cm}
\dirtree{%
  .1 Application.
  .2 Native Shell Layer.
  .3 Shell Container \emph{(top-level navigation state, lifecycle hooks)}.
  .3 Platform Bridge \emph{(explicit plugin allow-list, native facades)}.
  .2 Presentation Layer.
  .3 WebView Host \emph{(full-screen embedded browser bridge)}.
  .3 Blade Templates \& Assets \emph{(mobile layouts, Vite CSS/JS modules)}.
  .2 HTTP Boundary Layer.
  .3 Runtime Normalization \emph{(bootstrap fixes, body recovery, CSRF policy)}.
  .3 Policy Gate \emph{(cross-cutting access and compliance middleware)}.
  .3 Routing \emph{(explicit endpoints and catch-all fallback)}.
  .2 Application \& Service Layer.
  .3 Thin Controllers \emph{(request/response orchestration)}.
  .3 Stateless Services \emph{(transformations, helpers, static settings)}.
  .2 Data Access \& Background Layer.
  .3 Relational Store \emph{(Eloquent models, sessions, queue tables)}.
  .3 Encrypted Storage \emph{(portable key-value secrets)}.
  .3 Background Workers \emph{(generation-guarded queued jobs)}.
}
\hspace{1cm}

\noindent Translating these logical layers onto Laravel's skeleton gives every artifact an
unambiguous home. Native UI state lives in \q{app/NativeComponents/},
cross-cutting HTTP rules in \q{bootstrap/app.php} and \q{app/Http/Middleware/}, orchestration in
\q{app/Http/Controllers/}, stateless logic in \q{app/Services/} and \q{app/Helpers/}, static
configuration in \q{app/Settings/}, persistence in \q{app/Models/} and
\q{database/migrations/}, and asynchronous work in \q{app/Jobs/} and \q{app/Console/Commands/}.


\subsubsection{Building Generic Foundations}

Before implementing individual features, it pays to extract a small set of generic primitives
that every screen can rely on. Just as encapsulating UI inputs behind custom widgets prevents
"UX Hell" when libraries change\footnote{This mirrors how Object-Relational Mapping (ORM)
layers standardize database access: callers depend on one stable interface while the underlying
engine or storage mechanism can evolve independently.}, encapsulating HTTP helpers, fallback
routing, local secrets, and background loops keeps the rest of the codebase concise.

\paragraph{Base Controller and Catch-All Fallback}
When migrating or adapting a large multi-screen workflow, building every route before the first
release is rarely practical. Instead, we can define shared request primitives on an abstract
base \q{Controller} and terminate \q{routes/web.php} with a catch-all \q{GenericController}:

\begin{lstlisting}[language=php8]
// app/Http/Controllers/GenericController.php
final class GenericController extends Controller
{
  public function index(string $url)
  {
    $html = $this->fetch($url, request()->all());
    // ... Normalize relative assets and
    // strip pop-up targets for mobile WebView
    return view('generic', ['data' => $html]);
  }
}
\end{lstlisting}

\noindent In \q{routes/web.php}, specialized feature routes are declared first, followed by the
wildcard fallback:

\begin{lstlisting}[language=php8]
// Explicit mobile-adapted routes go first...
Route::get('/settings',
  [SettingsController::class, 'index']);

// ...followed by the catch-all fallback
Route::get('/{any}',
  [GenericController::class, 'index'])
    ->where('any', '.*');
Route::post('/{any}',
  [GenericController::class, 'indexPost'])
    ->where('any', '.*');
\end{lstlisting}

\noindent This simple pattern lets the application grow incrementally: high-traffic screens get
tailored controllers, stateless services, and mobile Blade templates, while secondary routes
continue functioning through the fallback pipeline without hitting dead-end 404 pages.

\paragraph{Stateless Services and Encapsulated Storage}
Controllers should orchestrate HTTP flow---validating input, calling a service, and picking a
view or redirect---while leaving data transformations and persistence rules to dedicated
classes. Keeping service classes under \q{app/Services/} stateless (accepting inputs, returning
plain arrays or DTOs, and never issuing redirects) makes them trivial to reuse across
controllers and queued jobs.

Similarly, sensitive local values such as credentials, tokens, or parental PINs should pass
through a single storage wrapper rather than calling filesystem functions or platform bridges
throughout the app:

\begin{lstlisting}[language=php8]
// app/Services/SecureStorage.php
class SecureStorage
{
  public static string $storagePath =
    'secure_storage_';

  public static function set(string $key, string $value): void
  {
    file_put_contents(
      storage_path(self::$storagePath . $key),
      encrypt($value)
    );
  }

  public static function get(string $key): ?string
  {
    $path = storage_path(self::$storagePath . $key);
    if (!file_exists($path)) {
      return null;
    }
    $encrypted = file_get_contents($path);
    return $encrypted ? decrypt($encrypted) : null;
  }

  public static function delete(string $key): void
  {
    $path = storage_path(self::$storagePath . $key);
    if (file_exists($path)) {
      unlink($path);
    }
  }
}
\end{lstlisting}

\noindent Two subtle design choices make \q{SecureStorage} especially effective here. First,
relying on Laravel's \q{encrypt()}/\q{decrypt()} (keyed by \q{APP\_KEY}) provides identical
behavior on Android, iOS, desktop \q{php artisan serve}, and automated tests. Second, exposing
\q{\$storagePath} as a configurable static prefix allows tests to redirect writes to an isolated
prefix during \q{setUp()} and clean them up in \q{tearDown()}.

\paragraph{Generation-Guarded Background Jobs}
Mobile apps frequently need long-lived or repeating background tasks (polling updates, syncing
local state, or listening to streams). Running them inside an HTTP request blocks the WebView,
whereas dispatching raw overlapping jobs risks duplicate workers after re-authentication. A
clean generic pattern guards a self-rescheduling \q{ShouldQueue} job with a random generation
token in cache:

\begin{lstlisting}[language=php8]
// app/Jobs/SyncWorker.php
class SyncWorker implements ShouldQueue
{
  private const CACHE_KEY = 'sync_worker_generation';

  public function __construct(public string $generation = '') {}

  public static function start(): void
  {
    $generation = bin2hex(random_bytes(8));
    Cache::forever(self::CACHE_KEY, $generation);
    static::dispatch($generation);
  }

  public function handle(): void
  {
    if (Cache::get(self::CACHE_KEY) !== $this->generation) {
      return; // Superseded by a newer chain or stopped
    }
    // ... perform incremental background work ...
    if (Cache::get(self::CACHE_KEY) === $this->generation) {
      static::dispatch($this->generation)->delay(now()->addSeconds(5));
    }
  }
}
\end{lstlisting}

\noindent Whenever \q{start()} is called again (or \q{Cache::forget(self::CACHE\_KEY)} stops
the worker), any previously queued instance exits cleanly on its next wake-up without needing
platform-specific job-cancellation APIs.


\subsubsection{Conceptualizing}

Because NativePHP embeds PHP directly inside a mobile host rather than behind Nginx and
PHP-FPM, several runtime boundaries require deliberate architectural decisions.

\paragraph{Single-Process Shell Guard}
In our hybrid shell setup (\cref{native-components}), \q{HomeShell} renders a full-screen
\q{<native:webview php>}. Both the native shell component and the pages loaded inside its
WebView execute inside the same PHP process. If a route inside the WebView redirects back to the
root entry screen and re-mounts \q{/shell/home}, the nested shell load can abort the runtime
(\q{SIGABRT}). To prevent this, \q{HomeShell::mount()} caches \q{getmypid()} under
\q{HomeShell::SHELL\_PID\_KEY}, and entry controllers check that process ID before deciding
whether to launch the shell or stay inside the active WebView:

\begin{lstlisting}[language=php8]
$shellExists = Cache::get(HomeShell::SHELL_PID_KEY) === getmypid();
return redirect($shellExists ? '/home' : '/shell/home');
\end{lstlisting}

\paragraph{Centralized HTTP Boundary Hotfixes and Policy Gates}
Mobile WebView bridges occasionally diverge from standard browser HTTP semantics. For instance,
NativePHP's iOS scheme handler can forward a \q{POST} body without setting a \q{Content-Length}
header, leaving \q{\$\_POST} empty, and WebView \q{fetch()} calls in NativePHP 4.3 do not
reliably attach session cookies. Rather than littering controllers with \q{php://input} parsing,
these quirks are isolated at the entry boundary: an early guard in \q{bootstrap/app.php}
populates \q{\$\_POST} from raw input, \q{HotfixForPostAction} middleware merges URL-encoded or
JSON bodies into Laravel's \q{Request}, and \q{validateCsrfTokens(except: ['*'])} is documented
in one place.

Cross-cutting access policies follow the same rule. Store-mandated checks such as neutral age
assurance or parental locks live in dedicated middleware (\q{EnsureAgeAssurance}) configured with
explicit \q{EXEMPT} and restricted route lists, ensuring no controller can accidentally bypass
compliance.

To separate intentional architecture, let's compare the implementation
with other applications listed under the
\href{https://github.com/topics/nativephp-mobile}{https://github.com/topics/nativephp-mobile}.
The official mobile starter
(\href{https://github.com/NativePHP/mobile-starter}{https://github.com/NativePHP/mobile-starter})
is a deliberately small Laravel baseline and is useful as a compatibility reference: it keeps
NativePHP and standard Laravel tooling preconfigured, automates dependency updates, and
publishes a ready iOS simulator build. The SuperStack starter
(\href{https://github.com/NativePHP/superstack}{https://github.com/NativePHP/superstack})
makes a more opinionated choice: Filament provides administration, \q{Route::native}
identifies SuperNative screens, and \q{nativephp/mobile-ui} supplies native UI on mobile while
the web UI renders the same Blade interface in a browser. The password manager example
(\href{https://github.com/codingwithrk/password-manager}{https://github.com/codingwithrk/password-manager})
demonstrates another valid hybrid: Laravel, Livewire, FluxUI, Tailwind, Alpine.js, and SQLite
are combined with paid native plugins for biometrics and secure storage.

\begin{table}[!ht]
  \centering
  \scriptsize
  \begin{tabularx}{\textwidth}{|p{0.22\textwidth}|X|X|}
    \hline
    \textbf{Architecture decision} & \textbf{Pros / cons of the alternative pattern} &
    \textbf{Our implementation and trade-off} \\
    \hline
    \textbf{ADR-1: Hybrid shell plus WebView} &
    A WebView-heavy UI is fast to port and shares PHP/Blade code with web. SuperStack follows this
    direction with native routes and browser fallback. It is less costly than building every view
    in Swift/Kotlin, but nested navigation and JavaScript-to-native state can become difficult. &
    \q{NavigationLayout} owns a small native tab surface and \q{HomeShell} loads Blade routes in
    an embedded WebView. This is ideal for a ported desktop-like game, but native-only animation,
    gesture fidelity, and platform-specific detail are limited. \\
    \hline
    \textbf{ADR-2: Bespoke routes first, generic fallback second} &
    Requiring an explicit controller and Blade view for every upstream route is safer and more
    predictable, but it makes initial migration prohibitively large. A fallback-only application
    is quick to build, yet hides layout and security regressions under a generic view. &
    \q{routes/web.php} declares adapted game workflows explicitly and sends unknown routes through
    \q{GenericController}. Pros: incremental delivery and no 404 dead ends; cons: fallback markup
    can silently diverge from the desktop site and legacy routes need periodic conversion. \\
    \hline
    \textbf{ADR-3: Laravel storage with a compatibility wrapper} &
    Device keychain/Keystore plugins provide platform-native secret isolation, as shown by the
    password-manager example, and biometrics can protect access. Their cons are paid or
    platform-specific packages, an additional registration step, and emulator or account
    requirements. & \q{SecureStorage} uses Laravel encryption over files in \q{storage/}. The
    implementation is free, portable, and testable across Android, iOS, and desktop. Its cons are
    that it depends on \q{APP\_KEY}, key rotation needs care, and it is an application wrapper
    rather than OS-backed hardware-protected storage. \\
    \hline
    \textbf{ADR-4: Explicit native plugin allow-list} &
    Auto-registering plugins reduces setup and matches conventional Laravel package discovery,
    which is convenient for starters. It can also increase the bundle and native attack surface
    because unrelated transitive packages become native capabilities. &
    \q{NativeServiceProvider::plugins()} explicitly enables the UI and Screen providers. This
    reduces surprise native dependencies and makes permission/capability reviews visible, while
    imposing a maintenance step whenever a new plugin is added. \\
    \hline
    \textbf{ADR-5: Generation-guarded Laravel queue jobs} &
    Using platform-native background APIs gives strong scheduling and lifecycle integration, but
    cancellation, permissions, and implementation details vary by OS. A raw Laravel queued loop
    is portable but cannot cancel an already queued instance by itself. &
    \q{ListenStream} uses a random cache generation, database-backed queueing, and delayed
    re-dispatch. Pros: restart/re-authentication invalidates stale chains, and desktop tests do not
    need Android jobs; cons: queue service lifecycle is host-dependent and 5-second rescheduling
    increases scheduling jitter. \\
    \hline
    \textbf{ADR-6: Laravel-first testing and web parity} &
    Device-only UI and native integration tests provide the highest fidelity but require simulators,
    certificates, and slow builds. Livewire/Flux and SuperNative approaches reduce client
    duplication differently, but both still need platform smoke tests. &
    Stateless services, overridable storage, native facade mocks, and PHPUnit cover most
    transformations and policies in seconds. Pros: fast feedback and CI-friendly regressions;
    cons: WebView HTTP, shell process, WebView cookies, and actual OS permissions are only proven
    after an APK or IPA run. \\
    \hline
  \end{tabularx}
  \caption{ADR-style comparison of architecture choices}
  \label{tbl:architecture-adrs}
\end{table}

The cross-project comparison suggests that none of the examined applications uses a monolithic
``native UI everywhere'' policy. The official starter optimizes for a lean baseline, SuperStack
optimizes for convention-led native and admin surfaces, and the password-manager example optimizes
for secure native capabilities. Our implementation follows the same Laravel-first principle, but
adapts it to an existing remote game: native code is restricted to navigation and platform
facades, while the bulk of feature behavior remains Blade, Laravel, SQLite, and stateless PHP.
This ADR is not a claim that all routes should remain in the generic fallback forever. The
boundary should move toward bespoke screens as a route becomes frequently used, performance
sensitive, security sensitive, or unsuitable for WebView behavior.

Putting these decisions into day-to-day practice comes down to six concrete guidelines:

\begin{enumerate}
  \item \textbf{Keep the native surface minimal.} Reach for \q{NativeComponent}s and device
  facades only when OS integration matters (tab bar, \q{Screen::keepAwake()},
  \q{Dialog::toast()}, \q{Device::vibrate()}). Keep forms, lists, and styling in Blade and Vite
  so they run under \q{php artisan serve} as well as on a phone.

  \item \textbf{Absorb platform quirks at the boundary.} Put scheme-handler fixes, encoding
  normalizers, and access gates in \q{bootstrap/app.php} and middleware. Never copy-paste
  \q{php://input} workarounds or permission checks into individual controllers.

  \item \textbf{Keep services stateless and controllers thin.} Services should accept primitive
  data or arrays and return structured results without touching \q{request()}, \q{session()}, or
  \q{redirect()}. Centralize static lookup tables and external URLs in \q{app/Settings/}.

  \item \textbf{Partition local state by lifecycle and sensitivity.} Use SQLite migrations for
  relational records, sessions (\q{SESSION\_DRIVER=database}), and queued jobs
  (\q{QUEUE\_CONNECTION=database}); use \q{SecureStorage} for secrets and compliance flags; and
  reserve \q{Cache} for transient process or worker coordination tokens.

  \item \textbf{Treat \q{nativephp/} as disposable.} Never rely on unrecorded edits inside the
  generated Android or iOS folders. Capture required host fixes in \q{\_patches/}, commit the
  PHP binary lock file, and register plugins explicitly in \q{NativeServiceProvider::plugins()}.

  \item \textbf{Verify behavior without a device first.} Because services are stateless, native
  facades can be mocked, and \q{SecureStorage::\$storagePath} is overridable, most regressions
  can be caught in seconds via PHPUnit before compiling an APK or IPA:
\end{enumerate}

\begin{lstlisting}[language=php8]
// tests/Unit/NativeComponents/HomeShellTest.php
public function test_pressing_the_active_tab(): void
{
  $screen = new HomeShell;
  $screen->showProfile();
  $first = $screen->webviewSrc();

  $screen->showProfile();
  $this->assertNotSame($first, $screen->webviewSrc());
}

// Feature test with mocked native facade
Dialog::shouldReceive('toast')->once();
Notification::addMessage('Saved successfully');
$this->assertDatabaseHas('notifications',
  ['message' => 'Saved successfully']);
\end{lstlisting}


